\chapter{Objetivos e motivação}


Este trabalho tem como objetivos apresentar:

\begin{itemize}
    \item Uma novam metodologia para a simulação de filmes finos, bem como as ferramentas e teorias utilizadas durante seu desenvolvimento.
    \item Os modelos matemáticos e numéricos utilizados, dando enfoque naquilo que de fato foi modificado ou implementado no código utilizado.
    \item Uma visão geral das teorias que tangenciam o método, mas que são necessárias para o entendimento do mesmo, justificativa da escolha de modelos ou interpretação dos resultados.
    \item Resultados comparativos alcançados com o novo método, identificando seu intervalo de validade, suas limitações e vantagens em relação aos demais disponíveis.
\end{itemize}

Nesta perspectiva, os capítulos dessa dissertação foram desenvolvidos com os seguintes objetivos. 

\begin{itemize}
    \item Revisão Bibliográfica: apresentar os trabalhos correlatos, bem como os principais problemas abordados pelos mesmos, assim como as teorias complementares necessárias para o entendimento das escolhas.
    \item 
\end{itemize}

O trabalho utilizou o código MFSim, desenvolvido pelo MFLab (UFU). Sendo assim, sua motivação surgiu da necessidade de resolver problemas envolvendo filmes finos em geometrias complexas, nos quais os métodos disponíveis no MFSim não apresentaram resultados satisfatórios. Além da motivação pela necessidade, também se almeja ter o novo método, como uma opção mais abrangente disponível para tratar de filmes finos. 